\begin{textsample}{NEG Dim 3 – global\_north\_2024\_12 – Score -1.00 – global\_north\_2024\_12/118334678.txt}  \label{ex:f3_neg_225}
A picture taken during a tour organised by the Israeli Army shows an Israeli truck transporting aid destined for the Gaza Strip at a drop-off area near the Kerem Shalom crossing, on November 28, 2024.
Jack Guez/AFP/Getty Images CNN -- At least 200 people were killed in Israeli airstrikes on northern Gaza Saturday, according to local health officials, as the United Nations said it would pause aid deliveries through the enclave ' s main crossing after more of its trucks were stolen.
The developments underscore the worsening humanitarian situation in the enclave, where tens of thousands of people have been killed by the Israeli military, and chronic hunger threatens the remaining civilian population.
On Friday, two children and a woman were crushed to death while attempting to buy food from a bakery in central Gaza.
The deadly strikes also come with an uneasy truce underway between Israel and Hezbollah in Lebanon, which Israeli Prime \textbf{Minister} Benjamin Netanyahu said would allow his forces to focus on Gaza.
Speaking to CNN on Sunday, Dr.
Hussam Abu Safiya, director of the Kamal Adwan Hospital, said that five buildings housing more than 200 people were struck in the Tal Al Zaatar and Beit Lahiya areas of northern Gaza."They were calling for help, and anyone who tried to assist was bombed.
Unfortunately, the cries for help have disappeared ; they were killed,"Dr.
Abu Safiya said.
The strike in Tel Al Zaatar left more than 100 people under the rubble, with only one person pulled out."This scene has become a daily occurrence, and no one is held accountable ; no one can stop the killing of innocent people."The spokesperson for Gaza ' s Civil Defense said that more than 40 people belonging to the"Al-Araj"family were killed in a single strike on a building in the Tel Al Zaatar neighborhood.
CNN has reached out to the Israel Defense Forces for further comment on the target of the strikes and the measures taken to mitigate civilian casualties.
According to the Health Ministry in Gaza, at least 44,429 people have been killed and more than 105,000 injured in the enclave since the war began last year.
The figure is thought to be an underestimate, as much of northern Gaza is inaccessible and many casualties never arrive at a hospital to be counted.
UN pauses aid deliveries The deadly strikes coincided with the theft of trucks carrying food and other supplies into the besieged strip, prompting the UN agency for Palestinian refugees to halt aid deliveries through the main crossing point between Israel and Gaza.
The"difficult decision"to stop deliveries through Kerem Shalom comes at a time when"hunger is rapidly deepening,"UNRWA chief Philippe Lazzarini warned Sunday.
The decision was made after a"few food trucks"were"taken"on the route Saturday, he wrote on X.
A source involved in transferring aid inside Gaza told CNN that a further five trucks loaded with flour were stolen near the crossing Sunday."The road out of this crossing has not been safe for months,"Lazzarini noted in his post, referring to an incident on November 16 when almost 100 aid trucks were stolen by armed gangs in what UNRWA described as"one of the worst"incidents of its kind.
The humanitarian operation in Gaza had become"unnecessarily impossible,"he added, citing hurdles from Israeli authorities and political decisions to restrict the amounts of aid as compounding factors in the breakdown of law and order in the enclave.
Lazzarini stressed that Israel, as the occupying power, was responsible for the protection of aid workers and supplies.
Israeli authorities"must ensure aid flows into Gaza safely and must refrain from attacks on humanitarian workers,"he said.
COGAT, the Israeli agency responsible for approving aid into Gaza, said that"dozens"of other humanitarian organizations continued delivering supplies to people in the enclave."Last week, over 1,000 trucks carrying humanitarian aid were collected from the various crossings and distributed throughout the Gaza Strip,"the agency added in a statement shared with CNN."We will continue to work with the international community to increase the amount of aid making its way into Gaza, through the Kerem Shalom Crossing as well as the other four crossings between Israel and Gaza."

% matched lemmas: minister
\end{textsample}
